\section{Preface to Persian Traditions}

First let me introduce myself. People who know me very well say that I have a mentality which is medieval and not modern, rural and not urban, that I am an ``incurable romantic and idealist'', and that I have a ``peasant mindset''. To all of the above I plead guilty, and as the Spanish say, \emph{y a mucha honra}, in other words, I am proud of it.

When one speaks of relations between Spain and Persia or of Persian influences in Spain, most people immediately think of elements which entered the Iberian Peninsula during the long period of Muslim dominance. This theme has been subject of many studies, by Hussein Munis and E. Levi-Provençal among others. For the above reasons, in the present study I am devoting much space to relations or influences which either predate the Muslim Conquest of Spain, which entered independently of said conquest, or, though they may have first entered with the Crescent of Islam, remained long after said Crescent had waned and set. To do a really complete and thorough study of this topic would require a great deal more time and money than I at present have at my disposal. In particular, it would require a long journey through Portugal, Aragon, Catalunya and Valencia, not to mention Iran. I have chosen to devote much time to Toledo, which was capital of Spain in Visigothic times and which is really a synthesis of all cultures, religions and artistic styles which form the threads of the multicolored fabric, part Celtic tartan, part Oriental carpet, which is the history of Spain.

I have also concentrated to a great degree on Asturias and Galicia. These regions were perhaps the least affected of all by the Muslim Conquest. Indeed, were it not for historical records there would be nothing to indicate that said conquest even took place in this Northwest corner of Spain. Yet, in the fields of art and architecture Sassanian influences were very strong indeed. Old Castile and Leon, lacerated for so long by endless border wars and border raids (of the men of Leon and Old Castile one might say what Sir Walter Scott said of the Vikings: ``They had a breastplate for a cradle and were fed from a blade.'') were unable to develop their own artistic and architectural styles. The art and architecture of Leon and Old Castile in the Middle Ages is simply one aspect of the European Romanesque and Gothic styles, lightly touched by Mozarab, Muslim and Mudejar influences. In the field of literature, I have concentrated on epic poetry simply because in the field of lyric verse it is practically impossible to determine what may be of Persian derivation and what is of Celtic and/or Provencal inspiration. Nevertheless, as we shall see, I have not completely neglected the lyric.

The Arabic language has no epic tradition, and therefore neither the Arabs nor their language could be considered as likely intermediaries between the epic traditions of Iran and Old Castile. I hope to demonstrate that one is really on quite firm ground when speaking of Iranian elements in the Castilian epic tradition. Although here one encounters the problem of how to distinguish between what is the result of Iranian influences and what proceeds from the strong Celtic substratum of Old Castile, I believe that it can be shown that there are a number of elements which the Castilian Epic shares only and exclusively with the Persian Epic.

I wish to thank all my friends in Galicia, Asturias, Andalusia, Extremadura, Old Castile and Leon for all the help and kindness which I received, and also to thank the Illustrious Mozarabic Community of Toledo for the kindness and hospitality which I received during the 1st International Congress of Mozarabic Studies, during which I was able to obtain photos of Toledo and its environs which help to illustrate this study. I also wish to thank the Southwestern College of Business of Middletown, Ohio and the Middletown Branch of Miami University (Ohio) for their kindness in allowing me to use their personal computers.

I also wish to dedicate this book to the memory of the late Walter Havighurst, my professor of creative writing at the University of Miami of Ohio. It was Professor Havighurst who really taught me the art of composition in English. In particular, it was Professor Havighurst's favorite motto, \textbf{\emph{Do not tell, show,}} which has guided me in my compositions, both fiction and non-fiction, as I recall several times in the present work.


\flrightit{Posted on 2020-03-13 by Michael McClain}

\begin{center}* * *\end{center}

\begin{footnotesize}
\begin{sffamily}
\texttt{Michael on 2020-03-17 at 09:37 said:}

Looking forward to this, thank you for publishing this online. Will attempt to actively follow along.

\hfill

\texttt{Rostam on 2020-03-27 at 23:44 said:}

Amazing! As an Iranian who is well read in European history, I was completely unaware of such influence. Looking forward to reading this book.


\hfill

\end{sffamily}
\end{footnotesize}

